% Intended LaTeX compiler: pdflatex
\documentclass[12pt]{article}

\usepackage[utf8]{inputenc}
\usepackage[T1]{fontenc}
\usepackage{graphicx}
\usepackage{longtable}
\usepackage{wrapfig}
\usepackage{rotating}
\usepackage[normalem]{ulem}
\usepackage{amsmath}
\usepackage{amssymb}
\usepackage{capt-of}
\usepackage{hyperref}
\usepackage{url}
\usepackage{sbc-template}
\usepackage[english]{babel}
\sloppy
\date{}
\title{Performance Analysis of Fletcher’s Reverse Time Migration Using Task-Based Paradigm on Multiple CPUs}
\hypersetup{
 pdfauthor={},
 pdftitle={},
 pdfkeywords={},
 pdfsubject={},
 pdfcreator={},
 pdflang={English}}
\begin{document}

\makeatletter
\let\orgtitle\@title
\makeatother

\title{\orgtitle}

\author{
   Pedro Henrique Boniatti Colle\inst{1},
   Lucas Mello Schnorr\inst{1},
   }

\address{
   Institute of Informatics, Federal University of Rio Grande do Sul (UFRGS)\\
   Caixa Postal 15.064 -- 91.501-970 -- Porto Alegre -- RS -- Brazil
   \email{schnorr@inf.ufrgs.br}}

\maketitle
\begin{abstract}
Put the abstract here.
\end{abstract}

\begin{resumo}
Coloque o resumo aqui.
\end{resumo}
\section{Introduction}
\label{sec:orgd8f5e6a}

This is just an example to show how Orgmode \cite{orgmode} can be
neatly used to write papers following the SBC Article latex style. Feel
free to modify and distributed this as you wish.

For a detailed explanation on why and how to write reproducible
articles using OrgMode, please visit the following. 
Check this paper \cite{schnorr2013visualizing} about how you need to

Sem subseções na introdução (max 2 página)

\begin{itemize}
\item Contextualização ->

Aplicação RTM é normalmente OpenMP parar explorar arq. multi-core
\begin{itemize}
\item Mas faz 15 anos que temos o paradgima task based (cita artigo don garra)
\end{itemize}

\item Não tem nenhuma solução hoje do fletcher em task based$\backslash$*  (tem uma outra implementação, já foi explorado)
\begin{itemize}
\item Por que o fletcher especificamente (Spadoto) -> sem task based
\begin{itemize}
\item Related work aqui
\end{itemize}
\item Eu fiz isso em CPU single node 
\begin{itemize}
\item Ainda é muito relevante fazer CPU-only (tá, como?) -> computação hybrida
\end{itemize}
\end{itemize}

Contribuições
\begin{itemize}
\item O programa
\item Otimização do IO
\item Análise de desempenho -> com múltiplas GPUS
\item Comparação com OpenMP base (que )
\end{itemize}

Estrutura do texto
\end{itemize}
\section{Your Great Contribution to Computer Science}
\label{sec:org85fd14b}
\label{org0eef2d5}
\section{Methods and Materials}
\label{sec:orgb22a0ba}

Figura do workflow do projeto\footnote{coisas2}
\subsection{Fletcher's Equations}
\label{sec:org95bd2dd}
\begin{itemize}
\item Equações diferenciais do fletcher.
\end{itemize}
\subsection{Task Based Implementation}
\label{sec:orgc502457}
\begin{itemize}
\item Descrição da implementação.
\begin{itemize}
\item Tudo o que eu fiz
\end{itemize}
\end{itemize}
\subsection{Methodology}
\label{sec:org4da4706}
\begin{itemize}
\item Projeto experimental
\begin{itemize}
\item Metodologia
\item máquinas
\item Nix e tal
\end{itemize}
\end{itemize}
\section{Results}
\label{sec:org76bffc0}
\label{orge4c024e}
\subsection{MSamples/s Performance}
\label{sec:orge1bb026}


\begin{center}
\includegraphics[width=.9\linewidth]{/tmp/nix-shell.63IGGu/babel-wLK0xp/figure8U04FG.pdf}
\label{org40e2210}
\end{center}
\subsection{Size increasse}
\label{sec:org3010ad2}

Gráficos do aumento do tamanho pelos Msamples

\begin{center}
\includegraphics[width=.9\linewidth]{/tmp/nix-shell.63IGGu/babel-wLK0xp/figure3HfoK7.pdf}
\label{org687e802}
\end{center}

Gráficos do aumento do tamanho pelo tempo de execução

\begin{center}
\includegraphics[width=.9\linewidth]{/tmp/nix-shell.63IGGu/babel-wLK0xp/figure6kTLXR.pdf}
\label{org190eeec}
\end{center}
\section{Conclusion}
\label{sec:org1a89bbf}

It was cool and awesome
\section*{Acknowledgments}
\label{sec:org744dbd8}
Who paid for this? stuff
\bibliographystyle{sbc}
\bibliography{refs}
\end{document}
